% ============================================================
% RIHB-in-the-loop, v4.
% Ground-up derivation: Fresnel coefficients -> body PO ->
% UE-anchored Kirchhoff render -> UL-pilot-calibrated
% closed loop -> pose-differentiable rate.
%
% Self-contained: no required reading of other docs.
% Author: R. Wydaeghe
% ============================================================

\documentclass[11pt,a4paper]{article}

\usepackage[margin=2.4cm]{geometry}
\usepackage[utf8]{inputenc}
\usepackage[T1]{fontenc}
\usepackage{lmodern}
\usepackage{microtype}
\linespread{1.04}

\usepackage{amsmath,amssymb,amsthm,mathtools}
\usepackage{bm}
\usepackage{mathrsfs}
\usepackage{bbm}
\usepackage{cancel}
\usepackage{xfrac}
\usepackage{booktabs}
\usepackage{multirow}
\allowdisplaybreaks[2]

\usepackage{empheq}
\usepackage[most]{tcolorbox}
\tcbset{highlight math style={%
    enhanced, colframe=black!70, colback=gray!6,
    arc=2pt, boxrule=0.5pt, left=3pt, right=3pt}}

\usepackage{siunitx}
\sisetup{detect-all, range-phrase = {--}, range-units = single,
         per-mode = symbol, inter-unit-product = {\,}}

\usepackage{enumitem}
\usepackage{xcolor}
\usepackage{xspace}
\usepackage{fancyhdr}
\pagestyle{fancy}
\fancyhf{}
\cfoot{\thepage}
\renewcommand{\headrulewidth}{0pt}

\usepackage{titlesec}
\titleformat{\section}{\large\bfseries}{\thesection}{1em}{}
\titleformat{\subsection}{\normalsize\bfseries}{\thesubsection}{1em}{}

\usepackage{tikz}
\usetikzlibrary{positioning,arrows.meta,calc,fit,backgrounds,shapes.geometric}

\usepackage[colorlinks=true, linkcolor={blue!55!black},
            citecolor={green!45!black}, urlcolor={blue!65!black}]{hyperref}
\usepackage[nameinlink,capitalize]{cleveref}

\usepackage{aliascnt}
\theoremstyle{definition}
\newtheorem{theorem}{Theorem}[section]
\newaliascnt{proposition}{theorem}
\newtheorem{proposition}[proposition]{Proposition}
\aliascntresetthe{proposition}
\newaliascnt{lemma}{theorem}
\newtheorem{lemma}[lemma]{Lemma}
\aliascntresetthe{lemma}
\newaliascnt{corollary}{theorem}
\newtheorem{corollary}[corollary]{Corollary}
\aliascntresetthe{corollary}
\newaliascnt{definition}{theorem}
\newtheorem{definition}[definition]{Definition}
\aliascntresetthe{definition}
\newaliascnt{remark}{theorem}
\newtheorem{remark}[remark]{Remark}
\aliascntresetthe{remark}
\crefname{proposition}{Proposition}{Propositions}
\Crefname{proposition}{Proposition}{Propositions}
\crefname{lemma}{Lemma}{Lemmas}\Crefname{lemma}{Lemma}{Lemmas}
\crefname{corollary}{Corollary}{Corollaries}\Crefname{corollary}{Corollary}{Corollaries}
\crefname{definition}{Definition}{Definitions}\Crefname{definition}{Definition}{Definitions}
\crefname{remark}{Remark}{Remarks}\Crefname{remark}{Remark}{Remarks}

\newcommand{\khat}{\hat{\bm{k}}}
\newcommand{\nhat}{\hat{\bm{n}}}
\newcommand{\rhat}{\hat{\bm{r}}}
\newcommand{\etahat}{\hat{\bm{\eta}}}
\newcommand{\rr}{\mathbf{r}}
\newcommand{\EE}{\mathbf{E}}
\newcommand{\HH}{\mathbf{H}}
\newcommand{\Sinc}{S_{\mathrm{inc}}}
\newcommand{\Sab}{S_{\mathrm{ab}}}
\DeclareMathOperator{\ReLU}{ReLU}
\DeclareMathOperator{\RE}{Re}
\DeclareMathOperator{\IM}{Im}
\DeclareMathOperator{\tr}{tr}
\DeclareMathOperator*{\argmin}{arg\,min}
\DeclareMathOperator*{\argmax}{arg\,max}
\newcommand{\diff}{\mathrm{d}}
\newcommand{\ident}{\mathbbm{1}}
\newcommand{\Reals}{\mathbb{R}}
\newcommand{\Complex}{\mathbb{C}}
\newcommand{\Qabs}{\mathbf{Q}_{\mathrm{ab}}}
\newcommand{\Phimat}{\mathbf{\Phi}}
\newcommand{\Mmat}{\mathbf{M}}
\newcommand{\Jmat}{\mathbf{J}}
\newcommand{\hh}{\mathbf{h}}
\newcommand{\xvec}{\mathbf{x}}
\newcommand{\thetab}{\bm{\theta}}
\newcommand{\betab}{\bm{\beta}}
\newcommand{\gammab}{\bm{\gamma}}
\newcommand{\Pabs}{P_{\mathrm{ab}}}
\newcommand{\SINR}{\mathrm{SINR}}
\newcommand{\Sigmaplus}{\Sigma_{+}}
\newcommand{\hbody}{h_{\mathrm{b}}}
\newcommand{\hLOS}{h_{\mathrm{LOS}}}
\newcommand{\Hkern}{\mathcal{K}}
\newcommand{\Tepskin}{T_{\mathrm{0}}}
\newcommand{\Frib}{\mathbf{F}_{\mathrm{b}}}
\newcommand{\Qref}{\mathbf{Q}_{\mathrm{re}}}
\newcommand{\kout}{\hat{\bm{k}}_{\mathrm{o}}}

\title{\textbf{The User as a Reflective Intelligent Body:\\
SAR receipts and pose-differentiable mmWave\\
capacity, from the Fresnel coefficient out (v4)}}
\author{Robin Wydaeghe\\
 Ghent University \& imec\\
 \texttt{robin.wydaeghe@ugent.be}}
\date{\today}

\makeatletter
\renewcommand{\maketitle}{%
  \thispagestyle{plain}%
  \begin{center}
    \rule{\linewidth}{1pt}\\[0.4em]
    {\LARGE\bfseries \@title}\\[0.5em]
    {\large \@author}\\[0.3em]
    {\small \@date}\\[0.4em]\rule{\linewidth}{1pt}
  \end{center}
  \vspace{0.5em}
}
\makeatother

\begin{document}
\maketitle

\begin{abstract}
A user holding a phone in an mmWave cell is, geometrically, a
reflective intelligent body: a near-array passive scatterer whose
surface phase pattern is fixed by pose, and whose only tuning
controller is the human brain.  This paper derives, from the
Fresnel reflection coefficient at a planar lossy interface, two
deployable products that share one piece of on-device machinery: a
live SAR receipt that displays the user's current absorbed power
in real time, and a pose-differentiable cascaded-channel rate
optimisation that suggests sub-degree to few-degree pose moves to
the user.  The derivation is self-contained.  Section~II builds
the per-triangle Fresnel transmission and reflection coefficients
under the pseudo-Brewster collapse that holds for skin tissue at
mmWave; Section~III lifts those per-triangle coefficients to a
discretised SMPL-X body surface; Section~IV writes the absorbed
power as a Hermitian quadratic form in the BS precoder; Section~V
derives the body-to-phone channel coefficient as a UE-anchored
Kirchhoff radiation integral, with no assumption that the body is
in any single far-field direction from the phone; Sections VI--VII
state the closed loop where uplink pilots calibrate the on-device
twin via a regional or modal regression on the body-mediated
channel residual; Sections VIII--IX develop the optimisation, the
compute architecture, and the bandwidth budget.  An empirical PoC
in the line-of-sight un-shadowing regime measures
\SI{14.5}{\percent} of rate gain at \SI{15}{\degree} of torso
yaw with no operationally relevant exposure cost.
\end{abstract}

\vspace{0.3em}
\noindent\textit{Style note.}~Wandering style, theorems where the
proof is short, propositions when a sketch is enough, conjectures
flagged.  Self-contained: every quantity is defined where it
appears, including the Fresnel coefficients in
\cref{sec:fresnel}.  Intended deployable home is \S\,IV-V of a
JSAC2 SI submission.

\tableofcontents
\bigskip

% ============================================================
\section{Introduction}\label{sec:intro}
% ============================================================

\subsection{The two reasons mmWave deployment slowed}

Two things slowed mmWave 5G past its original schedule.  First, the
human body in the line-of-sight of an FR2 channel costs
\SI{25}{} to \SI{40}{dB} of received power
\cite{MacCartney2017,Maccartney2017blockage}.  Bodies are
everywhere bodies are: torsos in offices, commuters in transit,
limbs near the held device.  Second, the post-deployment
public-relations conversation has been dominated by exposure
questions, and the conventional ``the regulator already certifies
it''~\cite{ICNIRP2020} response has not won the
argument.~\cite{Roosli2010, Roosli2021}.  Both deployment blockers
route through the human body in the cell.

This paper proposes that the same on-device body twin can address
both.  Part~1 turns the conversation by giving each user a
real-time, transparent dose receipt: \emph{this is exactly how
much electromagnetic power your body is currently absorbing.}
Part~2 turns the blockage cost into a small adaptable
channel-shaping opportunity: the body-as-passive-RIS, with pose
as the tuning knob the user already controls.  Both products use
the same per-triangle Fresnel-PO machinery; the differences are in
what is computed at each surface point and where the result is
displayed.

\subsection{Why ``reflective intelligent body''?}

A reconfigurable intelligent surface (RIS) is a passive panel of
sub-wavelength tunable phase shifters that reshapes incident waves
toward a desired receiver~\cite{DiRenzo2020,Wu2019}.  A human body
at mmWave is, in the high-frequency Kirchhoff/PO regime that
\cref{sec:po} formalises, a passive scatterer whose surface phase
pattern $e^{-ik_0\khat\cdot\rr}$ is \emph{fixed by pose} and whose
only tuning is the user's posture.  The body cannot tune
electronically; it can re-pose, on the order of \SI{1}{\Hz}.  In
one sentence:

\begin{quote}
\emph{Human bodies are reflectors at mmWaves.  They are also
mounted with a highly intelligent organ called the brain.  They
are therefore reflective intelligent human bodies (RIHBs).}
\end{quote}

We adopt RIS lingo throughout: ``cascaded channel,''
``passive RIS,'' ``surface phase mask,'' ``re-pointing.''  The
cascaded structure familiar from RIS,
\(
  \hh_{\mathrm{cascaded}} = \hh_{\mathrm{BS\to RIS}}\cdot
  \Phimat\cdot\hh_{\mathrm{RIS\to UE}},
\)
has a body-side analog
\begin{equation}\label{eq:cascaded-body}
  \hh(\thetab,\rr_p)=\hLOS(\rr_p)+\hbody(\thetab,\rr_p),
\end{equation}
where $\hbody$ is the body-mediated channel coefficient
\cref{sec:body-to-phone} derives below.  The body's
``$\Phimat$'' is the implicit per-triangle phase mask induced by
its surface and pose; the BS's ``re-pointing'' of the body is the
pose suggestion the user takes (or does not take).  The analogy
is informal but pedagogically useful and used throughout.

\subsection{Two contributions, one piece of machinery}

\begin{table}[h]
\caption{The two contributions of this paper, on the same per-triangle
Fresnel-PO machinery.}
\label{tab:two-parts}
\centering\small
\begin{tabular}{@{}p{0.30\linewidth}p{0.31\linewidth}p{0.31\linewidth}@{}}
\toprule
\textbf{Question} & \textbf{Part~1: SAR receipt} &
\textbf{Part~2: RIHB capacity} \\
\midrule
What is computed? &
$\Pabs(\thetab,\xvec)$ on the body surface in real time &
$\nabla_{\thetab} R(\thetab,\xvec)$ for a pose suggestion $\thetab^*$ \\
Output & a number, displayed in an app & a pose move, suggested via
notification \\
Where it lives & on the body & at the phone antenna \\
Privacy & receipt computed locally; opt-in aggregation &
gradient computed locally; only $\thetab^*$ accept/reject leaves
device \\
Optimization role & none: informational only &
maximises rate under comfort cost only \\
\bottomrule
\end{tabular}
\end{table}

% ============================================================
\section{Fresnel coefficients at a lossy half-space}\label{sec:fresnel}
% ============================================================

This section is the foundation everything else stands on.  We
derive the per-polarisation amplitude and power transmission and
reflection coefficients at a planar interface between vacuum and a
lossy dielectric, recover the pseudo-Brewster collapse that holds
for skin at mmWave, and write down the per-point absorbed-flux
identity that drives both the SAR receipt and the body-to-phone
channel.

\subsection{Setup}

Let region $z<0$ be vacuum and region $z>0$ be a lossy dielectric
with relative permittivity
$\tilde\varepsilon_r=\varepsilon_r-i\sigma/(\omega\varepsilon_0)$,
relative permeability $\mu_r=1$, and complex refractive index
$\tilde n=\sqrt{\tilde\varepsilon_r}$ (principal branch with
$\IM\tilde n\le 0$).  At \SI{28}{\GHz}, skin tissue has
$\tilde\varepsilon_r\approx 16.5-i12.6$ \cite{IT'IS2024} so
$\tilde n\approx 4.4-i1.4$.  An incident plane wave at the
interface is parametrised by polar angle $\theta_i$ from the
outward surface normal $\nhat=-\hat z$, with propagation direction
$\khat$ in the half-space $\khat\cdot\nhat<0$ and unit
amplitudes for the transverse-electric (TE / $s$) and
transverse-magnetic (TM / $p$) components.  Snell's law gives
$\sin\theta_i=\tilde n\sin\tilde\theta_t$ with complex transmitted
angle $\tilde\theta_t$.

\subsection{Amplitude coefficients}

Direct application of the boundary conditions $E_\parallel$
continuous and $H_\parallel$ continuous to the plane wave gives the
Fresnel amplitude transmission and reflection coefficients (TE in
the first column, TM in the second):
\begin{empheq}[box=\tcbhighmath]{equation}\label{eq:fresnel-amp}
\begin{aligned}
  t_s &= \frac{2\mu}{\mu+\xi},
  & r_s &= \frac{\mu-\xi}{\mu+\xi}, \\[2pt]
  t_p &= \frac{2\,\tilde n\,\mu}{\tilde n^2\,\mu+\xi},
  & r_p &= \frac{\tilde n^2\,\mu-\xi}{\tilde n^2\,\mu+\xi}.
\end{aligned}
\end{empheq}
where $\mu=\cos\theta_i$ and $\xi=\sqrt{\tilde n^2-1+\mu^2}$ with
$\RE\xi>0$.  These are exact for a planar interface and a lossy
half-space; no Born approximation is invoked.  The amplitudes are
linked by $t_s=1+r_s$ and $t_p=(1+r_p)/\tilde n$, which is
boundary-condition continuity in disguise.

\subsection{Power transmission and the absorbed flux}

The boundary-conditioned Poynting flux normal to the interface
gives, per polarisation, the power transmission
\begin{equation}\label{eq:T-power}
  T_s = \RE\!\left[\frac{\xi}{\mu}\right]\cdot |t_s|^2,\qquad
  T_p = \RE\!\left[\frac{\xi}{\tilde n^2\,\mu}\right]\cdot |t_p|^2.
\end{equation}
Power conservation per polarisation
$T_s+|r_s|^2=1$ and $T_p+|r_p|^2=1$ holds exactly, with the
absorbed flux per unit area at the interface
\begin{equation}\label{eq:Sab-perpol}
  S_\mathrm{ab}^{(\sigma)}(\theta_i) = \Sinc^{(\sigma)}\,T_\sigma(\theta_i)\,
  \cos\theta_i,
  \qquad \sigma\in\{s,p\},
\end{equation}
where $\Sinc^{(\sigma)}=|E_\sigma^{\mathrm{inc}}|^2/(2 Z_0)$ is the
incident Poynting flux of polarisation $\sigma$ in vacuum.  The
key observation for what follows is that $T_s$ and $T_p$ in
\cref{eq:T-power} depend on incidence angle and on tissue
permittivity but not on the precoder; they are deterministic
per-triangle scalars once the geometry is set.

\subsection{The pseudo-Brewster collapse}\label{sec:pBrewster}

\begin{lemma}[Pseudo-Brewster]\label{lem:pBrewster}
For a planar lossy half-space with refractive index
$|\tilde n|\gg 1$ (mmWave high-water tissue: $|\tilde n|\sim 5$ at
\SI{28}{\GHz}), and for incidence angles $\theta_i\in[0,80^\circ]$,
the polarisation-averaged power transmission is
\begin{empheq}[box=\tcbhighmath]{equation}\label{eq:pBrewster-T}
  T_s(\theta_i)\approx T_p(\theta_i)\approx
  T_0(\theta_i;\tilde n)
  \;=\; 1-\bigl|r_0(\theta_i;\tilde n)\bigr|^2
\end{empheq}
to within a relative error
$|T_s-T_p|/T_0\lesssim 5\,\%$ across the band, where
$|r_0|=\frac{1}{2}(|r_s|+|r_p|)\approx
|\tilde n-1|/|\tilde n+1|$ at small $\theta_i$.
At normal incidence and for skin at \SI{28}{\GHz},
$T_0(0)\approx 0.54$, hence $|r_0(0)|^2\approx 0.46$.  The angular
dependence is mild: $T_0$ stays in $[0.49,0.56]$ across
$\theta_i\in[0,70^\circ]$ before falling off near grazing.
\end{lemma}

The proof is direct from \cref{eq:fresnel-amp} and large-$|\tilde
n|$ algebra.  The practical consequence is that polarisation
sensitivity drops out of dosimetry at mmWave: each surface point
contributes the same scalar $T_0(\theta_i)$ regardless of
incident polarisation, and we can write
\begin{equation}\label{eq:Sab-pB}
  \Sab(\rr) = \Sinc(\rr)\,T_0\bigl(\theta_i(\rr)\bigr)
  \,\cos\theta_i(\rr)
  \quad
  \text{with}
  \quad
  \cos\theta_i(\rr)=-\khat\cdot\nhat(\rr).
\end{equation}
The reflected amplitude collapses likewise: in the
pseudo-Brewster regime,
\begin{equation}\label{eq:r-pB}
  r_s\approx r_p\approx r_0(\theta_i)\,e^{i\phi_r(\theta_i)}
  \quad \text{with}\quad |r_0|^2 = 1-T_0,
\end{equation}
with the angle-dependent phase $\phi_r$ acting as a per-triangle
scalar in what follows.

\subsection{What this section gives us}

Two scalars per surface point, computable from the surface normal
and the incident path direction alone, fully describe the local
EM book-keeping under the pseudo-Brewster regime:
\begin{itemize}[leftmargin=1.6em,itemsep=2pt]
\item the \emph{absorption fraction} $T_0(\theta_i)\cdot\cos\theta_i$,
which converts incident flux into absorbed flux pointwise
(\cref{eq:Sab-pB});
\item the \emph{reflected amplitude}
$\sqrt{1-T_0(\theta_i)}\,e^{i\phi_r(\theta_i)}$, which seeds the
Huygens secondary source that radiates body-to-UE
(\cref{eq:r-pB}).
\end{itemize}
Both are deterministic functions of geometry and tissue, with no
precoder dependence.  Sections \ref{sec:po}-\ref{sec:body-to-phone}
build everything else on top.

% ============================================================
\section{The body as a discretised PO scatterer}\label{sec:po}
% ============================================================

\subsection{The SMPL-X parametric body}

The body is modeled as a closed parametric surface
$\Sigma=\Sigma(\thetab,\betab)$ from SMPL-X
\cite{Pavlakos2019}.  The shape vector
$\betab\in\Reals^{10}$ parametrises bone lengths, body
proportions, and body-mass distribution; the pose vector
$\thetab\in\Reals^{72}$ parametrises 24 articulated joints in axis-angle
representation.  Linear blend skinning (LBS) maps a canonical
T-pose template mesh
$\Sigma_0\subset\Reals^3$ with $V=10\,475$ vertices and
$M_\triangle\sim 20\,000$ triangles to the posed surface
\begin{equation}\label{eq:lbs}
  \rr_v(\thetab,\betab)
  =\sum_{j=1}^{24} w_{vj}\,T_j(\thetab,\betab)\,
   \rr_v^{(0)}(\betab),
\end{equation}
with skinning weights $\{w_{vj}\}$ supplied by SMPL-X and per-joint
rigid transforms $T_j(\thetab,\betab)\in SE(3)$ that compose
$24\times 4\times 4$ floats per pose.  The smooth dependence
$\thetab\mapsto\rr_v$ is $C^\infty$ on the open set of
non-singular kinematics and is the differentiability foundation
\cref{sec:diff} relies on.

\paragraph{Per-triangle local frame.}  Each triangle
$t\in\{1,\ldots,M_\triangle\}$ has vertices
$\{\rr_{v_1(t)},\rr_{v_2(t)},\rr_{v_3(t)}\}$, centroid
$\rr_t=\tfrac{1}{3}\sum_i\rr_{v_i(t)}$, area $A_t$, outward unit
normal $\nhat_t$, and an associated tangent pair
$(\hat e_{1,t},\hat e_{2,t})$.  All four are smooth functions of
$(\thetab,\betab)$ via \cref{eq:lbs}.

\paragraph{Visibility from a point.}  The PO regime requires that
the surface element be lit by the source: triangles in the body's
self-shadow must not contribute.  We define the visibility
indicator
\begin{equation}\label{eq:vis-from}
  V(\rr;\rr_s)=\begin{cases}
    1 & \text{the segment from $\rr_s$ to $\rr$ is not occluded by $\Sigma$},\\
    0 & \text{otherwise.}
  \end{cases}
\end{equation}
$V(\rr;\rr_s)$ is computed by ray-mesh intersection
($O(\log M_\triangle)$ per query with a BVH).  It is non-smooth at
silhouette transitions; we replace the indicator by a per-triangle
GELU smoothed gate $V_\delta$ with width $\delta\sim 1^\circ$ in
incidence angle, giving a $C^\infty$ surrogate.  The bias on
$\Pabs$ from the gate softening is bounded by $\delta\cdot
\partial(\nhat\cdot(-\khat))/\partial\theta\cdot
\mathrm{(silhouette\ length)}$ and stays under \SI{0.5}{\percent}
empirically.

\subsection{Single-bounce PO, with the corrections we ignore}

We adopt single-bounce physical optics for the body's interaction
with the radiation: each lit triangle reflects according to the
local Fresnel coefficients (\cref{lem:pBrewster}) and is
absorbed according to \cref{eq:Sab-pB}.  Three contributions are
neglected:
\begin{itemize}[leftmargin=1.6em,itemsep=2pt]
\item evanescent and Norton-wave tangential surface fields.
Negligible at mmWave (skin depth \SI{0.7}{mm}; reactive zone
\SI{1.7}{mm}); reduces \cref{eq:Sab-pB} accuracy by less than
\SI{1}{\percent} except in the immediate-contact regime;
\item creeping and diffraction around limb contours: introduces a
smoothing of the visibility gate already absorbed into $V_\delta$
above;
\item multi-bounce body re-illumination (torso reflects to arm,
arm absorbs): a Neumann series in the body's albedo $1-T_0\sim
0.46$ gives a bounded correction $\le (1-T_0)\bar V\sim
\SI{15}{\percent}$ on the absorbed power and is left for
\cref{sec:threads}.
\end{itemize}

\subsection{The base-station array}

A single base station has $M$ antenna elements at known positions
$\{\rr_j\}_{j=1}^M$.  In running examples $M=64$ ($8\times 8$
uniform rectangular array, half-wavelength spacing, carrier
$f_c=\SI{28}{\GHz}$, $\lambda=\SI{1.07}{\cm}$, aperture
\SI{0.14}{\m}).  The element pattern $g_j(\hat\xi)$ is taken as a
calibrated complex vector function of direction with the per-element
gain factor absorbed in.  The array radiates with precoder
$\xvec\in\Complex^M$ subject to $\|\xvec\|^2\le P_{\mathrm{tx}}$.

\paragraph{The far-field assumption on the BS side.}  The Fraunhofer
distance for an aperture of diameter $D$ at carrier $f_c$ is
$r_F=2D^2/\lambda$.  For the BS aperture
($D=\SI{0.14}{\m}$, $\lambda=\SI{1.07}{\cm}$),
$r_F^{\mathrm{BS}}=\SI{3.7}{\m}$.  Bodies at the cell range
($\sim\SI{30}{\m}$) sit comfortably in the BS array's far field,
and the radiated waves arrive at the body as plane waves with a
single direction and a single polarisation amplitude per element,
modulated by the per-element pattern $g_j(\hat\xi)$.

\paragraph{The path dictionary.}  Propagation between the array
and a probe point $\rr$ is described by a list of
$N$ paths
\(
  \mathcal{D}(\rr)=\{j(n),\khat_n,\bm\psi_n\}_{n=1}^N,
\)
where $j(n)\in\{1,\ldots,M\}$ is the source element, $\khat_n$ is
the unit propagation direction at the probe, and $\bm\psi_n\in
\Complex^3$ is the complex polarisation amplitude (carrying any
intervening reflections, refractions, and the per-element
pattern).  The total incident field at the probe is the sum
\begin{equation}\label{eq:E-inc-probe}
  \EE^{\mathrm{inc}}(\rr;\xvec)
  =\sum_{n=1}^N x_{j(n)}\,\bm\psi_n(\rr)\,
   e^{-ik_0\khat_n\cdot\rr}.
\end{equation}
The dictionary $\mathcal{D}(\rr)$ is sampled at the probe location
by a ray tracer that traverses the wireless scene.  We treat
$\mathcal{D}$ as given for the rest of this paper; how the BS
shares it with the phone is the bandwidth budget of
\cref{sec:bandwidth}.

% ============================================================
\section{The absorbed power side: the Qabs operator}\label{sec:Qabs}
% ============================================================

This section converts the per-triangle absorption flux of
\cref{eq:Sab-pB} into a Hermitian quadratic form in the precoder.
The result is the operator $\Qabs(\thetab)$ that the SAR receipt
displays and that v5 (the BS-side compliance paper) projects
against.

\subsection{Per-triangle incident flux}

Substitute \cref{eq:E-inc-probe} into the definition of incident
Poynting flux at triangle $t$:
\begin{equation}\label{eq:Sinc-t}
  \Sinc(\rr_t;\xvec)
  =\frac{1}{2 Z_0}\bigl|\EE^{\mathrm{inc}}(\rr_t;\xvec)\bigr|^2
  =\frac{1}{2 Z_0}\sum_{n,n'} x_{j(n)}\overline{x_{j(n')}}\,
    \bm\psi_n^H\bm\psi_{n'}\,
    e^{-ik_0(\khat_n-\khat_{n'})\cdot\rr_t}.
\end{equation}
The sum has $N^2$ cross-terms, all of which carry phase factors
$e^{-ik_0(\khat_n-\khat_{n'})\cdot\rr_t}$ that depend on the
triangle position $\rr_t$.

\subsection{Per-triangle absorbed power}

Multiply \cref{eq:Sinc-t} by the per-triangle Fresnel weight from
\cref{eq:Sab-pB} and the triangle area:
\begin{equation}\label{eq:Pab-t}
  \Pabs^{(t)}(\xvec)
  =\Sab(\rr_t;\xvec)\cdot A_t
  =A_t\,T_0\!\bigl(\theta_t\bigr)\,
   \frac{\cos\theta_t^{(\mathrm{eff})}}{2 Z_0}
   \sum_{n,n'} x_{j(n)}\overline{x_{j(n')}}\,
    [\cdots].
\end{equation}
Here $\theta_t$ stands for the per-path incidence angle, which
\emph{varies per path} on the same triangle; in
\cref{eq:Sab-pB} we used a single $\theta_i$ for a single
incident wave.  The honest generalisation drops the
single-incidence simplification:
\begin{equation}\label{eq:Pab-t-fullcoupling}
  \Pabs^{(t)}(\xvec)
  =\frac{A_t}{2 Z_0}
   \sum_{n,n'} x_{j(n)}\overline{x_{j(n')}}\,
   \tilde T_{nn'}(\rr_t)\,
   \bm\psi_n^H\bm\psi_{n'}\,
   e^{-ik_0(\khat_n-\khat_{n'})\cdot\rr_t},
\end{equation}
with the per-path-pair Fresnel weight
\begin{equation}\label{eq:Tnnp}
  \tilde T_{nn'}(\rr_t)
  = \sqrt{T_0(\theta_n(\rr_t))\,T_0(\theta_{n'}(\rr_t))}
    \cdot \sqrt{\mu_n^+(\rr_t)\,\mu_{n'}^+(\rr_t)},
\end{equation}
$\mu_n^+ = \max(0,-\khat_n\cdot\nhat_t)$ the visibility-and-cosine
gate.

\subsection{Sum over the body}

Sum \cref{eq:Pab-t-fullcoupling} over the lit triangles and define
the per-body Hermitian Gram
\begin{empheq}[box=\tcbhighmath]{equation}\label{eq:Qabs-def}
  \Qabs(\thetab)\in\Complex^{M\times M},\quad
  [\Qabs]_{j,j'}\;=\;\sum_{n:j(n)=j}\sum_{n':j(n')=j'}\!\!
   M_{nn'}(\thetab),
\end{empheq}
with
\begin{equation}\label{eq:Mnnp}
  M_{nn'}(\thetab)
  =\frac{1}{2 Z_0}
   \int_\Sigma A(\rr)\,
   \tilde T_{nn'}(\rr)\,
   \bm\psi_n^H\bm\psi_{n'}\,
   e^{-ik_0(\khat_n-\khat_{n'})\cdot\rr}
   \diff A(\rr),
\end{equation}
where $A(\rr)$ is a continuous interpolation of the per-triangle
area density (one factor of $A_t/3$ at each vertex; vanishes
elsewhere).  The whole-body absorbed power is the Hermitian
quadratic form
\begin{empheq}[box=\tcbhighmath]{equation}\label{eq:Pabs-Qabs}
  \Pabs(\thetab,\xvec) = \xvec^H\,\Qabs(\thetab)\,\xvec.
\end{empheq}
The matrix $\Qabs$ is positive semidefinite by construction (a
Gram), $C^\infty$ in $\thetab$ except on the measure-zero set of
silhouette transitions (smoothed away by $V_\delta$), and has
effective rank set by the angular resolution of the BS array on
the body solid angle (typically rank 3-5 at the \SI{99}{\percent}
trace fraction; this is empirical).

\subsection{Two enabling approximations}\label{sec:approxs}

The Hermitian Gram structure of \cref{eq:Qabs-def} relies on two
non-trivial approximations that the absorbed-power literature has
identified
\cite{Wydaeghe2025monograph}.  They are stated here for
completeness; the error budget is bounded and small at mmWave.

\begin{proposition}[Approximation 1: TM polarisation direction]
\label{prop:approx1}
Inside a lossy half-space with refractive index $|\tilde n|\gg 1$,
the refracted TM direction $\hat e'_{p,n}$ differs from the
incident TM direction $\hat e_{p,n}$ by an angle
$\theta_n-\theta_{t,n}$ with $\sin\theta_{t,n}=
\sin\theta_n/|\tilde n|$.  The two normal components differ by
$\sin\theta_n(1-1/|\tilde n|)$ but the dot product
$\hat e'_{p,n}\cdot\hat e'_{p,n'}$ that appears in cross-terms is
dominated by the tangential product and the normal-normal product
is $O(1/|\tilde n|^2)\lesssim 4\,\%$ for skin at \SI{28}{\GHz}.  We
replace $\hat e'_{p,n}$ by $\hat e_{p,n}$.  Self-terms are
unaffected.
\end{proposition}

\begin{proposition}[Approximation 2: universal depth coupling]
\label{prop:approx2}
The cross-pair depth integral
\(
  \Lambda_{nn'}=\int_0^\infty e^{-(\alpha_n+\alpha_{n'})z}\,
  e^{i(\beta_n-\beta_{n'})z}\diff z
\)
that emerges from the ohmic-loss integral in the lossy half-space
factorises as $\Lambda_{nn'}=\sqrt{\Lambda_{nn}\Lambda_{n'n'}}\,
\Gamma_{nn'}$ with the normalised coupling
$\Gamma_{nn'}=2\sqrt{\alpha_n\alpha_{n'}}/(\alpha_n+\alpha_{n'}-
i(\beta_{n'}-\beta_n))$.  For high-water tissue at mmWave, the
spread in $\alpha_n$ across the path dictionary is small, and
$\Gamma_{nn'}\approx 1$ uniformly with relative error
$\le\SI{2}{\percent}$.  We absorb this into the per-path-pair
Fresnel weight $\tilde T_{nn'}$ of \cref{eq:Tnnp}.
\end{proposition}

The two approximations reduce the per-triangle absorbed flux from
a triple sum (path pair $\times$ polarisation direction $\times$
depth coupling) to a clean per-path-pair quadratic form, which is
what makes the operator $\Qabs$ a single Hermitian Gram rather
than a depth-coupled tensor.  Both errors are bounded and
non-conservative (they slightly under-estimate $\Pabs$); see
\cite{Wydaeghe2025monograph} for the full error budget.

\subsection{Translation phasor identity}\label{sec:translation}

The body moves through the cell.  Per-slot rebuild of
\cref{eq:Mnnp} is too expensive for a deployment: the surface
integral is $O(M_\triangle\cdot N^2)$.  We exploit one identity.

Let $\rr=\rr_0+\mathbf{t}$ where $\rr_0$ are the body-local mesh
coordinates (constant in absolute body translation) and
$\mathbf{t}$ is the body's world-frame translation per-slot.
Substituting into \cref{eq:Mnnp},
\begin{empheq}[box=\tcbhighmath]{equation}\label{eq:translation}
  M_{nn'}(\rr_0+\mathbf{t},\thetab)
  =e^{-ik_0(\khat_n-\khat_{n'})\cdot\mathbf{t}}\,
   M_{nn'}(\rr_0,\thetab),
\end{empheq}
which says $\Mmat$ is built once per pose at the body-local frame
($\sim$\SI{7.5}{\ms} per body on a phone-class GPU empirically), and
per-slot translation is an $O(N^2)$ phasor sandwich that costs
sub-millisecond.  This is what makes the SAR receipt deployable
at slot cadence.  An identical identity holds for the body-to-phone
channel of \cref{sec:body-to-phone} and gives that side the same
deployability.

\subsection{What the absorbed-power operator is and is not}

$\Qabs$ is the operator that maps every BS precoder to the body's
absorbed power, with the body geometry baked into the matrix.  It
is the SAR receipt's primary computation: $\Pabs=\xvec^H\Qabs\xvec$.
It is also v5's compliance constraint: the BS holds
$\Pabs\le L_{\mathrm{RL}}$ via the per-slot projection
$\xvec\to\alpha\xvec$ with $\alpha$ chosen to satisfy
\cref{eq:Pabs-Qabs}.

$\Qabs$ is \emph{not} the body-to-UE channel.  It counts power
\emph{landing in the body}, not power \emph{leaving the body
toward the phone}.  The body-to-UE machinery (which the rate
gradient of \cref{sec:diff} needs) is a different object,
derived next.

% ============================================================
\section{The body-to-phone channel: a UE-anchored Kirchhoff render}\label{sec:body-to-phone}
% ============================================================

This section is the new contribution relative to the back-hemisphere
$\Frib(\thetab,\kout)$ scattering matrix that prior work (and this
paper's earlier drafts) imported from radar.  We do not assume the
body re-radiates into a single far-field direction $\kout$.  We
compute the field at the phone antenna directly, summing over the
body's surface as a collection of Huygens secondary sources, and
weighting by the per-direction-of-arrival UE antenna pattern.

\subsection{The PO surface current}

For a single incident plane wave from direction $\khat_n$ with
polarisation amplitude $\bm\psi_n$ at a triangle $t$, the
specularly reflected amplitude in the upper half-space is, in the
pseudo-Brewster regime (\cref{eq:r-pB}),
\begin{equation}\label{eq:E-refl}
  \EE^{\mathrm{refl}}_n(\rr_t)
  = \sqrt{1-T_0(\theta_n^{(t)})}\,e^{i\phi_r(\theta_n^{(t)})}\,
    \mathrm{P}_n^{(t)}\,\bm\psi_n,
\end{equation}
where $\theta_n^{(t)}=\arccos(-\khat_n\cdot\nhat_t)$ is the local
incidence angle, and $\mathrm{P}_n^{(t)}$ is the per-polarisation
projector that maps the incident polarisation onto the local
(TE,TM) basis at the triangle and applies the corresponding
pseudo-Brewster amplitude.  In the pseudo-Brewster regime
$\mathrm{P}_n^{(t)}$ collapses to the planar projector
$(\ident_3-\khat_n^{\mathrm{r}}\khat_n^{\mathrm{r}\,T})$ where
$\khat_n^{\mathrm{r}}=\khat_n+2(\khat_n\cdot\nhat_t)\nhat_t$ is
the reflected direction.

The per-triangle equivalent surface current that supports the
reflected field, by Stratton-Chu / PO theory \cite{Balanis2005}, is
$\mathbf J_t=2\nhat_t\times\HH_n^{\mathrm{inc}}$ on the lit half
$\Sigma_+(\khat_n)$ and zero elsewhere.  We write the reflected
amplitude $\EE^{\mathrm{refl}}_n$ with the prefactor of
\cref{eq:E-refl} and treat each triangle as a point Huygens
source of dipole moment proportional to
$\EE^{\mathrm{refl}}_n(\rr_t)\,A_t$.

\subsection{Propagation to the phone}

The phone receiver is at world position $\rr_p$, with antenna
pattern $\bm g_{\mathrm{UE}}(\hat\xi)$ that returns the
polarisation-aware complex gain for a wave arriving from direction
$\hat\xi$.  The free-space scalar Green's function from
$\rr$ to $\rr_p$ is
$G(\rr,\rr_p)=e^{ik_0|\rr-\rr_p|}/(4\pi|\rr-\rr_p|)$.  The per-point
direction of arrival at the phone is
\begin{empheq}[box=\tcbhighmath]{equation}\label{eq:eta-perpoint}
  \etahat(\rr;\rr_p) \;\equiv\; \frac{\rr_p-\rr}{|\rr-\rr_p|}.
\end{empheq}
This direction \emph{varies pointwise on $\Sigma$}: each lit
triangle has its own direction-of-arrival at the phone, its own
distance, and its own angle into the UE antenna pattern.  This is
the central observation that distinguishes this section from the
$\Frib(\thetab,\kout)$-with-single-$\kout$ formulation.

\paragraph{The Kirchhoff radiation integral.}  By the
Stratton-Chu radiation formula \cite{Balanis2005,Senior1995}, the
field at $\rr_p$ from the PO surface current and reflected
amplitude on $\Sigma_+(\khat_n)$ is
\begin{equation}\label{eq:KirchhoffPO}
  \EE^{(n)}_{\mathrm{body}}(\rr_p)
  = ik_0\!\!\int_{\Sigma_+(\khat_n)}\!\!
   \bigl(\ident_3-\etahat\,\etahat^T\bigr)\,
   \EE^{\mathrm{refl}}_n(\rr)\,
   G(\rr,\rr_p)\,
   V(\rr;\rr_p)\,
   \diff A(\rr)
\end{equation}
where $V(\rr;\rr_p)$ is the body-self-occlusion gate from the phone
side (the source-side counterpart of \cref{eq:vis-from},
implemented by the same BVH-accelerated ray-mesh visibility query).
The integrand is the Huygens secondary source amplitude
$\EE^{\mathrm{refl}}_n$, projected to the transverse plane at the
arrival direction $\etahat$, propagated by the free-space Green
function, and gated by self-visibility from the phone.

\paragraph{The phone-side polarisation projection.}  The complex
voltage at the UE antenna terminals from the body-mediated wave is
the inner product of the incident field at the phone with the
phone's effective height vector
$\bm g_{\mathrm{UE}}(\etahat)$:
\begin{equation}\label{eq:UE-projection}
  v^{(n)}_{\mathrm{body}}(\rr_p)
  = \bm g_{\mathrm{UE}}\bigl(\etahat\bigr)^H \cdot
    \EE^{(n)}_{\mathrm{body}}(\rr_p).
\end{equation}
This is where the per-point $\etahat(\rr;\rr_p)$ matters
operationally: the UE antenna gain depends on the arrival direction,
and that direction varies across the body surface.  The integral
\cref{eq:KirchhoffPO} cannot be factored as
$(\text{single-direction gain})\times(\text{integral over body})$
when the body subtends a non-trivial solid angle from the phone, as
it does at smartphone-on-body geometries.

\subsection{The body-mediated channel coefficient}

Putting \cref{eq:KirchhoffPO} and \cref{eq:UE-projection} together
and summing over BS paths gives the per-element body-mediated
channel:
\begin{empheq}[box=\tcbhighmath]{equation}\label{eq:hbody}
  \boxed{\;
  \hbody^{(j)}(\thetab,\rr_p)
  =\!\!\sum_{n:j(n)=j}\!\!
   \int_{\Sigma_+(\khat_n)}
   K_n(\rr;\rr_p,\thetab)\,
   V(\rr;\rr_p)\,
   \bm\psi_n^H\,\bm g_{\mathrm{UE}}\!\bigl(\etahat\bigr)\,
   \frac{e^{ik_0|\rr-\rr_p|}}{4\pi|\rr-\rr_p|}\,
   e^{-ik_0\khat_n\cdot\rr}\,
   \diff A(\rr).
  \;}
\end{empheq}
The per-triangle scalar
\begin{equation}\label{eq:K-kernel}
  K_n(\rr;\rr_p,\thetab)
  = ik_0\,\sqrt{1-T_0\!\bigl(\theta_n^{(t)}\bigr)}\,
    e^{i\phi_r(\theta_n^{(t)})}\,
    \bigl(\hat e^{(\sigma)}_n\cdot\bigl(\ident-\etahat\,\etahat^T\bigr)
    \cdot\hat e^{(\sigma)}_n\bigr)
\end{equation}
absorbs the local Fresnel reflection amplitude, its angle-dependent
phase, and the polarisation projection onto the
$\etahat$-transverse plane.  In the pseudo-Brewster regime the
polarisation factor in $K_n$ is unity for the dominant
field-aligned component and the kernel reduces to
$ik_0\sqrt{1-T_0}\,e^{i\phi_r}$, scalar in the precoder index.

The per-element coefficients of \cref{eq:hbody} stack into
$\hbody(\thetab,\rr_p)\in\Complex^M$.  The cascaded channel of
\cref{eq:cascaded-body} is then
\begin{equation}\label{eq:hcascaded}
  \hh(\thetab,\rr_p) = \hLOS(\rr_p) + \hbody(\thetab,\rr_p)
   \;\in\;\Complex^M,
\end{equation}
with $\hLOS$ the direct (not body-mediated) channel coefficient,
itself a sum over the BS-to-phone paths in $\mathcal{D}(\rr_p)$
that do not interact with the user's body.  In the strict-LOS
case $\hLOS$ is one path; in NLOS it is the sum of all wall and
ceiling reflections that reach the phone without the body's
surface being a node.

\subsection{What the integral does and does not assume}

\Cref{eq:hbody} is computed without invoking any far-field
assumption on the body-to-phone propagation.  The integrand is
exact PO; the only approximations are:
\begin{itemize}[leftmargin=1.6em,itemsep=2pt]
\item single-bounce body-to-UE (no body-to-body re-illumination);
\item pseudo-Brewster polarisation collapse, which sets
$|r_s|\approx|r_p|\approx\sqrt{1-T_0}$;
\item GELU-smoothed visibility gates (\cref{eq:vis-from});
\item locally-plane-wave incidence on each triangle from the BS,
which holds because the BS is in the body's far field
($r_F^{\mathrm{BS}}=\SI{3.7}{\m}\ll\SI{30}{\m}$).
\end{itemize}
None of these assumes the phone is in the body's far field or
introduces a single $\kout$ direction.  At the phone-on-body
geometry where the torso subtends $\sim\SI{40}{\degree}$ of solid
angle from the phone at \SI{30}{\cm}, every triangle has a
distinct $\etahat$ and a distinct $|\rr-\rr_p|$, and the integral
is computed with these per-triangle quantities.

\subsection{The same translation identity holds}

The translation identity of \cref{eq:translation} carries to
\cref{eq:hbody} verbatim: substituting
$\rr=\rr_0+\mathbf{t}$ and noting that $\etahat$ and
$|\rr-\rr_p|$ both depend only on the relative geometry
$\rr_0-\rr_p$ when the body translates rigidly,
\begin{equation}\label{eq:hbody-trans}
  \hbody^{(j)}(\rr_0+\mathbf{t};\thetab)
  = \!\!\sum_{n:j(n)=j}\!\!
    e^{-ik_0\khat_n\cdot\mathbf{t}}\;
    \hbody^{(j),(n)}(\rr_0;\thetab),
\end{equation}
where the per-path pre-built coefficient
$\hbody^{(j),(n)}(\rr_0;\thetab)$ is the integral over the
body-local mesh in body-local coordinates.  Per-slot translation
is again a per-path phasor multiplication, $O(N M)$, sub-millisecond
on phone-class hardware.

% ============================================================
\section{UE-anchored rendering: a complexity argument}\label{sec:ue-render}
% ============================================================

\subsection{Forward radiosity vs.\ backward path tracing}

Path tracing in computer graphics renders an image of a 3D scene
by tracing rays \emph{from the camera backward into the scene},
because most rays leaving the scene do not reach the camera and
the camera-anchored render culls them for free.  An equivalent
forward renderer that traces rays from light sources is
correct but operationally wasteful when the only quantity wanted
is the field at the camera (i.e., one pixel).

The same observation governs $\Qabs$ vs.\ $\hbody$:
\begin{itemize}[leftmargin=1.6em,itemsep=2pt]
\item \textbf{$\Qabs(\thetab)$ is BS-anchored.}  Computed once per
body, it is an $M\times M$ Gram of the surface integral, used by
the BS to enforce compliance against \emph{every precoder it might
radiate}.  Cost per body refresh:
$O(N^2\cdot M_\triangle)$ for the integral plus $O(M^2)$ for the
matrix view.  The BS amortises this across many bodies and many
slots; deployment-feasible.
\item \textbf{$\hbody(\thetab,\rr_p)$ is UE-anchored.}  Computed
on the phone, it is one complex coefficient per BS element per
body slot.  The integral
\cref{eq:hbody} is summed over only the triangles visible from the
phone (visibility cull from $\rr_p$).  Per-element cost:
$O(M_\triangle^{(\mathrm{vis})}\cdot N^{(\mathrm{vis})})$, where
$M_\triangle^{(\mathrm{vis})}$ is the visible triangle count from
$\rr_p$ and $N^{(\mathrm{vis})}$ is the number of BS paths that
illuminate at least one of those visible triangles.
\end{itemize}

\subsection{Empirical cost on a phone-on-body geometry}

For a typical phone-held configuration (phone at \SI{30}{\cm} from
torso centroid, oriented toward face), the visibility cull from
$\rr_p$ removes about \SI{80}{\percent} of the body mesh
($M_\triangle^{(\mathrm{vis})}\sim\num{4000}$ on a
\num{20000}-triangle SMPL-X mesh).  The number of BS paths
illuminating at least one of those triangles depends on the
deployment: at the typical mmWave LOS-dominated geometry,
$N^{(\mathrm{vis})}\le\num{20}$ of $N=50$ total paths.  So the
phone-side cost per BS element is
\begin{align*}
  \text{per element}: \quad &
  M_\triangle^{(\mathrm{vis})}\cdot N^{(\mathrm{vis})}\sim\num{8e4}
  \text{ complex multiplies}\\
  \text{per slot}: \quad & M\cdot M_\triangle^{(\mathrm{vis})}\cdot N^{(\mathrm{vis})}
  \sim \num{5e6} \text{ complex multiplies}.
\end{align*}
At \SI{32}{bit} complex precision and a \SI{200}{GFLOPS} phone-class
NPU/GPU \cite{Apple2024A18}, the per-slot cost is
$\sim$\SI{50}{\micro\second}, well below the \SI{1}{\milli\second}
PHY frame and the \SI{30}{\milli\second} slot cadence.  The
visibility query (BVH ray-mesh intersection from $\rr_p$ for each
of the BS paths' arrival directions) is
$O(N\log M_\triangle)\sim\SI{1}{\milli\second}$ per slot.  Total
per-slot phone-side budget is $\sim$\SI{1.5}{\milli\second}, with
$\Mmat$-rebuild and pose-gradient steps amortised at
$\le\SI{1}{\Hz}$ refresh.

\subsection{The complexity-class point}

In one sentence: the BS-anchored $\Qabs$ render does
$O(M^2 M_\triangle)$ work per body to handle every precoder it
might use; the UE-anchored $\hbody$ render does
$O(M\cdot M_\triangle^{(\mathrm{vis})})$ work per body per slot
because \emph{the only field point that matters is the phone}.
This is what makes the on-device twin tractable.  It is also why
the radar-style $\Frib(\thetab,\kout)$ Gram of prior work is the
wrong primitive for the body-to-UE channel: $\Frib$ is built to
handle every $\kout\in S^2$, but the deployment only ever
queries one direction (the phone's).

% ============================================================
\section{The closed loop: UL pilots calibrate the twin}\label{sec:closedloop}
% ============================================================

\subsection{Why the loop needs closing}

\Cref{eq:hbody} gives the on-device prediction $\hbody^{\mathrm{pred}}(\thetab,\rr_p)$
from the BS-shared dictionary, the SMPL-X body, the UE pattern, and
the user's pose estimate.  Three sources of error blur this
prediction relative to the channel the user actually observes:
\begin{enumerate}[leftmargin=1.6em,itemsep=2pt]
\item \emph{BS-side dictionary amplitude error}: the ray tracer's
material codes and 3D-tile geometry have systematic errors that
miscount per-path amplitudes;
\item \emph{Body-side shape and surface error}: SMPL-X is a
generic parametric mesh; per-user surface roughness, clothing,
non-skin tissues introduce per-region amplitude and phase
deviations from the canonical SMPL-X surface;
\item \emph{Pose error}: the IMU-and-AR pose pipeline has finite
RMS accuracy ($\sigma_\theta\sim 4$--$8^\circ$ per joint).
\end{enumerate}
Without measurement, the on-device prediction is a model.  With
measurement, it becomes a calibrated estimator.  The measurement
is uplink pilots: the phone transmits known pilots, the BS
measures the cascaded channel
$\hh^{\mathrm{meas}}\in\Complex^M$ at $\sim\SI{1}{\milli\second}$
PHY-frame cadence, and the residual against the BS's own
prediction closes the loop.

\subsection{The residual decomposition}

Let $\hh^{\mathrm{pred}}(\thetab;\bm\beta,\bm\gamma)$ denote the
BS-side prediction of the cascaded channel under the same
machinery the phone runs, parametrised by:
\begin{itemize}[leftmargin=1.6em,itemsep=2pt]
\item $\bm\beta\in\Complex^N$: per-path complex amplitude
calibration on the BS dictionary
(this is the v5 \emph{dictionary calibration}, fit by ridge LS
against served-user uplink CSI as
$\hat{\bm\beta}=\argmin_{\bm\beta}\|\hh^{\mathrm{meas}}-
\mathrm{diag}(\bm\beta)\hh^{\mathrm{pred}}\|^2+\rho\|\bm\beta-\mathbf{1}\|^2$);
\item $\bm\gamma$: per-body calibration for the body-side terms,
introduced and stratified below.
\end{itemize}
The innovation
\(
  \Delta = \hh^{\mathrm{meas}} - \hh^{\mathrm{pred}}(\thetab;\bm\beta,\bm\gamma)
\)
splits, by linearity of \cref{eq:hcascaded}, into a non-body
contribution (already absorbed by $\bm\beta$ on the dictionary
side) and a body contribution that $\bm\gamma$ targets.

\subsection{Choosing the body-side parametrisation}\label{sec:gamma-choice}

We considered three parametrisations of $\bm\gamma$ and tested them
empirically (see \cref{sec:eval-cal} for the recovery curves).
The relevant trade-off is identifiability vs.\ expressiveness, and
the regime that controls it is the rank of
$\Lambda(\thetab,\rr_p)\in\Complex^{M\times M_\triangle^{(\mathrm{vis})}}$,
the per-element-per-triangle Kirchhoff coefficient matrix at the
current pose.

\paragraph{Three candidate parametrisations.}
\begin{enumerate}[leftmargin=1.6em,itemsep=2pt]
\item \emph{Single complex scalar} ($\gamma_0\in\Complex$).
$\hbody^{\mathrm{cal}}=\gamma_0\hbody^{\mathrm{pred}}$.  One DOF.
Identifiable from any UL pilot.  Captures only an overall amplitude
and phase miscalibration.
\item \emph{Per-region scaling} ($\gammab\in\Complex^{R}$, $R=6$
anatomical regions).  $\hbody^{\mathrm{cal}}=\sum_b\gamma_b
\hbody^{(b)}$ where $\hbody^{(b)}$ is the partial Kirchhoff sum
restricted to triangles in region $b$.  $R$ DOF.  Identifiable when
$\Lambda$ has rank $\ge R$; otherwise the regional bases collapse
and the LS becomes ill-conditioned.
\item \emph{$K$-mode SVD coefficients} ($\bm\gamma^{\mathrm{B}}\in\Complex^K$).
SVD $\Lambda=U\Sigma V^H$, truncate to the first $K$ modes:
$\hbody^{\mathrm{cal}}=\sum_{k=1}^K\gamma_k^{\mathrm{B}}
\sigma_k U_k$.  $K$ DOF.  The basis is orthogonal and adaptive to
the pose; identifiability is automatic when $K\le M$.
\end{enumerate}
For each, the fit is ridge LS against the residual
$\hh^{\mathrm{meas}}-\hLOS-\hbody^{\mathrm{pred}}$, with a small
ridge weight $\rho$ that biases toward the un-perturbed twin.

\paragraph{Empirical winner (\cref{sec:eval-cal}).}
The best parametrisation is geometry-dependent:

\begin{itemize}[leftmargin=1.6em,itemsep=2pt]
\item In the \emph{far-BS regime} (BS at \SI{30}{\m},
body-to-BS angular footprint sub-beamwidth, $K_{99}\!=\!1$),
$\Lambda$ is effectively rank one; the regional basis collapses
($\hbody^{(b)}\propto u_1$ for all $b$), and Tier C error plateaus
above Tier D.  The $K\!=\!8$ SVD parametrisation, in contrast,
reaches \SI{1}{\percent} body-channel reconstruction error at
\SI{40}{\dB} pilot SNR.
\item In the \emph{close-BS regime} (BS at \SI{5}{\m}, body
super-resolved, $K_{99}\!=\!4$), per-region $\bm\gamma\in\Complex^6$
is the most identifiable: the array sees regional differences
distinctly and Tier C reaches \SI{3}{\percent} at \SI{40}{\dB} SNR,
beating both the single scalar and the SVD parametrisation.
\end{itemize}

\paragraph{What we deploy.}  $K$-mode SVD with
$K=K_{99}(\thetab,\rr_p)$ chosen at each pose refresh tick from
the singular spectrum of $\Lambda$.  The empirical $K_{99}$ ranges
from 1 (far-BS, body sub-beamwidth) to 4 (close-BS, body
super-resolved); $K\le 10$ across all tested geometries.  The fit
is ridge LS at $\rho=0$ (the system is overdetermined since
$K\ll M$); the SVD itself is amortised across slots
($\sim$\SI{80}{\milli\second} per pose refresh on phone-class
hardware).  The DL bandwidth for refined $\bm\gamma^{\mathrm{B}}$
is $\le 10$ floats per body per second
$\le\SI{0.3}{\kilo bps}$, negligible relative to the path
dictionary.

\subsection{Why this is principled}

Three observations make $K$-mode SVD with
$K\!=\!K_{99}$ the principled choice rather than a hyperparameter
guess:
\begin{enumerate}[leftmargin=1.6em,itemsep=2pt]
\item \emph{Adaptivity.}  $K_{99}$ is computed from the singular
spectrum at the current pose, not chosen up front.  As the BS-body
geometry changes, $K$ tracks.
\item \emph{Identifiability automatic.}  $K\le M$ by construction;
the LS is overdetermined and never ill-conditioned, regardless of
$\Lambda$'s structural rank.
\item \emph{Optimality in the perturbation class.}  For a Tier C
perturbation $\bm\gamma^{(\mathrm{C})}$ in the regional basis, the
$K_{99}$-truncated SVD captures \SI{99}{\percent} of the
representable variation in $\hbody$.  The reconstruction error
floor is then determined by the pilot noise, not by the
parametrisation.
\end{enumerate}

\subsection{Why the loop is now closed}

The phone's prediction $\hbody^{\mathrm{pred}}(\thetab,\rr_p)$ is
periodically corrected by $\gammab$ values fit from
measurements the BS already collects.  The corrected estimate
$\hbody^{\mathrm{cal}}(\thetab,\rr_p;\hat\gammab)$ is the basis
for both the SAR receipt and the pose-gradient suggestion.  When
the user moves a joint by $\delta\thetab$, the change in cascaded
channel can be predicted in advance from the calibrated twin and
verified after the fact from the new UL pilot measurement; the
gradient computation rides on the calibrated estimate and is
self-validating.

\subsection{Pointer to v5's beta calibration}\label{sec:csi-cal-pointer}

The dictionary-side calibration $\bm\beta$ is identical in form to
v5's CSI calibration (v5~\S\,III.E): per-path complex scale fit
by ridge LS against served-user uplink CSI.  We do not re-derive
it here; we re-use it as the dictionary-side companion to the
$\bm\gamma$ ladder.  The $\bm\gamma$ side is the new contribution.
Both run at the BS at $\sim 1\,\mathrm{s}$ cadence; the path
dictionary, $\bm\beta$, and $\bm\gamma$ are all broadcast down to
the phone in the closed-loop control plane of
\cref{sec:bandwidth}.

% ============================================================
\section{Pose-differentiable rate}\label{sec:diff}
% ============================================================

\subsection{The cascaded SINR and the rate model}

For single-stream MRT precoding $\xvec=\sqrt{P_{\mathrm{tx}}}
\hh(\thetab)/\|\hh(\thetab)\|$ on the cascaded channel
$\hh(\thetab)$ of \cref{eq:hcascaded}, the SINR is
\begin{equation}\label{eq:SINR}
  \SINR(\thetab) = \frac{P_{\mathrm{tx}}\|\hh(\thetab)\|^2/\sigma_\eta^2}
  {1+\sum_{u\ne k}\text{interference at $u$}/\text{etc.}},
\end{equation}
which, in the deployable single-user single-stream regime, reduces
to $\SINR(\thetab)=P_{\mathrm{tx}}\|\hh(\thetab)\|^2/\sigma_\eta^2$.
The rate is the operational 5G NR FR2 rate
$R(\thetab)=B\cdot\min(\log_2(1+\SINR),\,S_{\max})$ with
$S_{\max}=\SI{7.4}{bps/Hz}$ (MCS27, \cite{TS38214}).  The cap matters
for the gradient interpretation: large SINR gains compress to
small rate gains at the cap, so the rate gradient saturates above
$\sim\SI{22}{dB}$ SINR.

\subsection{The pose gradient}

The cascaded channel $\hh(\thetab)$ is differentiable in pose by
the chain
\begin{equation}\label{eq:chain}
  \frac{\partial\hh}{\partial\theta_i}
  = \underbrace{\frac{\partial\hLOS}{\partial\theta_i}}_{\text{LOS un-shadowing}}
  + \underbrace{\frac{\partial\hbody}{\partial\theta_i}}_{\text{body-mediated}},
\end{equation}
both terms computed by autograd through the SMPL-X chain
$\thetab\to\{\rr_t(\thetab)\}\to\hbody$ and $\thetab\to v(\thetab)$
(visibility), respectively.  The rate gradient is
\begin{empheq}[box=\tcbhighmath]{equation}\label{eq:rate-grad}
  \frac{\partial R}{\partial\theta_i}
  = \frac{B}{\ln 2}\cdot
    \frac{P_{\mathrm{tx}}/\sigma_\eta^2}
         {1+P_{\mathrm{tx}}\|\hh\|^2/\sigma_\eta^2}\cdot
    2\,\RE\!\left(\hh^H\,\frac{\partial\hh}{\partial\theta_i}\right),
\end{empheq}
provided the operating SINR is below the MCS cap; above, the
gradient is zero in the cap-saturated regime.

\subsection{Differentiability statement}

\begin{proposition}[$C^\infty$ pose dependence]\label{prop:smooth}
The map $\thetab\mapsto(\hh(\thetab),\Qabs(\thetab))$ is
$C^\infty$ on the open set of poses where (i) no SMPL-X
joint is at a kinematic singularity and (ii) no surface triangle
crosses a visibility silhouette as $\thetab$ moves.  Under the
GELU smoothing of \cref{eq:vis-from} with width $\delta=1^\circ$,
condition (ii) is replaced by global $C^\infty$ at the cost of a
relative bias on $\Pabs$ and on $\hbody$ bounded by
$\delta\cdot(1-T_0)\cdot(\text{silhouette length}/\sqrt{A_{\mathrm{body}}})$,
$\le\SI{0.5}{\percent}$ in our PoC geometries.
\end{proposition}

The proof is direct from the SMPL-X $C^\infty$ statement
(\cref{eq:lbs}) and the integrability of \cref{eq:hbody} and
\cref{eq:Qabs-def} in pose, given the smoothed visibility gates.
The chain rule then assembles the gradient of any scalar functional
of $\hh$ and $\Qabs$ that is itself $C^\infty$, including
\cref{eq:rate-grad} and the SAR receipt $\Pabs$.

\subsection{The user-side optimisation loss}

The optimisation that drives the pose suggestion is
\begin{empheq}[box=\tcbhighmath]{equation}\label{eq:loss}
  \mathcal{L}(\thetab,\xvec;\thetab^{(0)})
  = -R(\thetab,\xvec)
   \;+\;\lambda_M\,C(\thetab-\thetab^{(0)}),
\end{empheq}
where $C$ is a separable comfort cost calibrated such that
$C\le 1$ corresponds to a $5^\circ$ torso rotation (RULA score
increment $+1$, ``low load'' \cite{McAtamney1993}), and
$\lambda_M$ trades pose-comfort against rate.  No exposure penalty
appears in \cref{eq:loss}.  In operational downlink, the user is
\emph{always} below the ICNIRP basic restriction
\cite{ICNIRP2020,FCCSAR2024}: typical handset-served-from-BS
configurations show $\Pabs\sim\SI{0.001}\times L_{\mathrm{RL}}$;
the receipt is informational, not a constraint that the
optimisation needs to enforce.  Compliance enforcement, if and
when it binds, is the BS-side projection of v5 (\cref{sec:saridx}).

\subsection{Where the gradient comes from}

The two terms of \cref{eq:chain} have different magnitudes in
different geometries.  In LOS-dominated body-shadowed scenarios
(the PoC of \cref{sec:eval-poc}), $\partial\hLOS/\partial\theta_i$
dominates: pose moves un-shadow the direct path, the cascaded
channel grows, the rate grows.  In strict NLOS scenarios where
$\hLOS=0$, $\partial\hbody/\partial\theta_i$ is the only term and
the gradient drives the user toward poses where the body
re-radiates a useful path to the phone.  In multi-user
co-operative scenarios (\cref{sec:eval-multi-user}), the gradient
on user $A$'s pose includes user $B$'s body in the cascaded
channel and vice versa; the joint optimum is a correlated
equilibrium.

% ============================================================
\section{The closed-loop block diagram}\label{sec:loop}
% ============================================================

\begin{center}
\begin{tikzpicture}[
  node distance=8mm and 14mm, font=\small,
  block/.style={draw, rounded corners, align=center, minimum height=8mm,
                minimum width=20mm, fill=gray!8},
  arrow/.style={-{Stealth[length=2mm]}, thick},
]
\node[block] (bs) {BS array};
\node[block, right=of bs] (rt) {Ray tracer};
\node[block, right=of rt] (dict) {Path dict.\\$\mathcal{D}$};
\node[block, below=of dict] (uephone) {Phone twin\\(SMPL-X + Kirchhoff)};
\node[block, below=of bs] (precoder) {Precoder\\(proj.\ ZF)};
\node[block, below=of precoder] (comp) {Compliance\\enforcer};
\node[block, right=of uephone, xshift=8mm] (loss) {Rate loss\\$+$ comfort};
\node[block, right=of loss] (suggest) {Pose\\suggest};
\node[block, below=of uephone] (sar) {SAR receipt\\display};
\node[block, below=of dict, yshift=-22mm] (cal) {$\bm\beta,\bm\gamma$\\calibrator};
\draw[arrow] (rt)--(dict);
\draw[arrow] (dict)--(uephone);
\draw[arrow] (uephone)--(loss);
\draw[arrow] (loss)--(suggest);
\draw[arrow] (suggest.south) -- ++(0,-7mm) -| (rt.south);
\draw[arrow] (uephone)--(sar);
\draw[arrow] (rt)--(precoder);
\draw[arrow] (precoder)--(comp);
\draw[arrow] (comp)--(bs);
\draw[arrow, dashed] (suggest) -- ++(0,12mm) -| (uephone.north);
\draw[arrow] (bs.east) -- ++(-2mm,-12mm) |- (cal.west);
\draw[arrow] (cal.east) -| (uephone.south);
\end{tikzpicture}
\end{center}

\noindent The two new blocks relative to v5 are the on-device twin
(SMPL-X + Kirchhoff render at the phone) and the
$\bm\beta,\bm\gamma$ calibrator.  Solid arrows are the data plane;
the dashed arrow is the control-plane pose suggestion the user
accepts or rejects.

\subsection{Bandwidth budget}\label{sec:bandwidth}

\begin{table}[h]
\centering\small
\caption{Bandwidth budget for the closed loop, per-user, per-second.
Path dictionary stays within the existing v5 control-plane
allocation; $\bm\beta,\bm\gamma$ are the new payload.}
\label{tab:bandwidth}
\begin{tabular}{@{}p{0.40\linewidth}p{0.27\linewidth}p{0.20\linewidth}@{}}
\toprule
Channel & Payload & Rate \\
\midrule
DL: path dictionary           & $50\times 8$ floats $\times \SI{10}{Hz}$ & \SI{128}{kbps} \\
DL: dictionary calibration $\bm\beta$ & $50$ floats $\times$ \SI{1}{Hz} & \SI{1.6}{kbps} \\
DL: body calibration $\bm\gamma$ (Tier C) & 6 floats $\times$ \SI{1}{Hz} & \SI{0.2}{kbps} \\
DL: body calibration $\bm\gamma$ (Tier B) & 20 floats $\times$ \SI{1}{Hz} & \SI{0.7}{kbps} \\
DL: pose suggestion $\thetab^*$ & $72$ floats $\times$ \SI{1}{Hz} & \SI{2.3}{kbps} \\
UL: pose estimate $\hat\thetab$ & $72$ floats $\times$ \SI{1}{Hz} & \SI{2.3}{kbps} \\
UL: pilots (already present)    & in PHY frame                       & 0 incremental \\
UL: accept/reject + SAR opt-in  & one bit + occasional bits          & negligible \\
UL: shape $\bm\beta_{\mathrm{body}}$ (one-time) & 10 floats once / session & \SI{40}{B} \\
\bottomrule
\end{tabular}
\end{table}

\noindent Total bidirectional control bandwidth $\le\SI{135}{kbps}$,
four orders of magnitude below the data rate the loop negotiates
over.

\subsection{Privacy delta against v5}

The phone holds the SMPL-X body and runs the Kirchhoff render
locally; only the pose-suggestion accept/reject bit and the (opt-in)
aggregated SAR receipt leave the device.  The BS already has the
pose estimate (UL), the path dictionary (it owns it), and the
calibrated $\bm\beta,\bm\gamma$ values (it computes them from UL
pilots).  The user's body geometry is implicit in $\bm\gamma$ but
is not directly disclosed: $\bm\gamma_b$ is a per-region complex
scale that the BS reads as a six-vector, not a mesh.  The privacy
delta over v5 is thus: same input data flow, plus on-device
pose-gradient compute, plus opt-in aggregated receipt.

% ============================================================
\section{Part 1: the live SAR receipt}\label{sec:saridx}
% ============================================================

\subsection{What it computes}

At each slot, the phone evaluates
\begin{equation}\label{eq:sar-receipt}
  \Pabs(\thetab,\xvec_t;\hat{\bm\gamma}_t)
  = \xvec_t^H\,\Qabs^{\mathrm{cal}}(\thetab;\hat{\bm\gamma}_t)\,\xvec_t
  \;[\mathrm{W}],
\end{equation}
where $\Qabs^{\mathrm{cal}}$ is the calibrated per-body operator
of \cref{sec:closedloop}: $\Qabs$ scaled by the appropriate
$\hat{\bm\gamma}$ on its body-region or modal decomposition, with
the same calibration flowing into the $\hbody$ side of the loop
in \cref{sec:diff}.  $\xvec_t$ is the precoder the BS used at the
slot, broadcast in the path-dictionary update plane.  The user
sees the receipt in absolute Watts, optionally normalised to (a)
the local reference-level budget $L_{\mathrm{RL}}$, (b) a body-mass
SAR $\Pabs/m_{\mathrm{body}}$, or (c) a 6-min time average for the
Brussels regulatory window.

\subsection{Why this has not been done}

Standard SAR disclosure on handsets is a fixed worst-case
$\mathrm{SAR}_{1g}$ or $\mathrm{SAR}_{10g}$ measured under
standardised conditions in a phantom.  Real-time per-user SAR
requires (i) real-time ray tracing of the user's local channel,
(ii) a calibrated body model on-device, and (iii) the BS sharing
the path information that real-time SAR depends on.  All three
become available when the BS already runs a body twin for
compliance: the BS is doing the ray tracing already, SMPL-X is on
the phone for AR purposes anyway, and the BS-to-UE control plane
already carries DL telemetry that can absorb the path dictionary
and calibration values.

\subsection{What the receipt is and what it is not}

The receipt is a model-based estimate of the user's absorbed
power, calibrated against the BS's own measurement of the
cascaded channel.  It is not a regulatory measurement; it does not
replace the basic-restriction-pathway compliance audit that the
regulator certifies.  It is a transparent number the user can act
on (move, re-orient the device, defer a session), and a rich data
source for population-scale dose epidemiology of the kind that
\cite{ICNIRP2020}~Annex~A~\S\,1.4 admits the field has been
unable to collect.

\subsection{Aggregated receipts and dose-effect epidemiology}

Per-user logged receipts, with consent, become input to dose-effect
studies of mmWave exposure that have been operationally impossible
without per-user, per-slot data.  The receipt aggregator is opt-in
and can be implemented with standard differential-privacy primitives
(noise injection on the per-slot sample; secure aggregation across
users).  The study-design question is its own paper.

% ============================================================
\section{App UX: gradient-with-feedback, not one-shot pose targeting}\label{sec:app}
% ============================================================

The on-device twin has two outputs: an absorbed-power scalar
(receipt) and a pose gradient (capacity coach).  This section
sketches the user-facing app and explains why differentiability
matters for the UX even when the global optimum could in principle
be precomputed.

\subsection{Why one-shot pose targeting fails}

In a perfect-twin idealisation, the phone could solve
$\thetab^*=\argmin_\thetab\mathcal{L}$ from
\eqref{eq:loss} once and instruct the user to assume $\thetab^*$
verbatim.  Three error sources break this in practice:
\begin{enumerate}[leftmargin=1.6em,itemsep=2pt]
\item \emph{Twin-model bias.}  SMPL-X is a canonical-mean
parametric body; per-user surface roughness, clothing dielectric,
and small-feature deviations from the canonical mesh produce a
per-region bias on $\hbody^{\mathrm{pred}}$ that we cannot measure
off-line.
\item \emph{Pose-estimator noise.}  The IMU + AR pose pipeline has
RMS error $\sigma_\theta\sim 4$--$8^\circ$/joint.  Even if the
predicted optimum is right, the achieved pose is
$\hat\thetab=\thetab^*+\bm\epsilon$.
\item \emph{Path-dictionary drift.}  As the user moves through the
cell, BS scatterers come and go; the dictionary's per-path
amplitudes are non-stationary.
\end{enumerate}
Compounded, the predicted optimum $\thetab^*_{\mathrm{pred}}$ is
significantly away from the actual maximum-rate pose
$\thetab^*_{\mathrm{true}}$.

\subsection{The closed loop converts ``find optimum'' into ``ride a
calibrated local model''}

The differentiable + UL-pilot-calibrated twin reframes
optimisation as a noisy-gradient descent on a slowly-drifting
non-stationary objective:
\begin{enumerate}[leftmargin=1.6em,itemsep=2pt]
\item Phone computes $\nabla_{\thetab}R$ from the current
calibrated twin.
\item App suggests a small move
$\delta\thetab=-\alpha\nabla_{\thetab}\mathcal{L}/\|\nabla\mathcal{L}\|$.
\item User attempts $\delta\thetab$; achieves
$\delta\hat\thetab=\delta\thetab+\bm\epsilon$.
\item UL pilots after the move give $\hh^{\mathrm{meas}}$;
calibration step refits $\bm\gamma$ against the residual.
\item Refined twin computes the next gradient; repeat.
\end{enumerate}
Convergence in 2--5 iterations is the empirical observation for
the geometries of \cref{sec:eval-poc}.  Note that step 4 is
\emph{free}: the UL pilots are in the PHY frame anyway; their
information is what the operator already has.

\subsection{The metal-detector metaphor}

The UX pattern is exactly that of a metal detector or a
PT-biofeedback fitness coach.  The user moves; the system responds
in real time with a haptic or visual ``warmer/colder'' signal
proportional to $\|\nabla R\|$ at the current pose.  The user does
not need to nail an exact target $\thetab^*$ (which they could not
reach anyway given $\sigma_\theta$); they need only know the
direction of improvement and follow it.  The twin's confidence in
the gradient direction grows as UL pilots accumulate, so early
moves are coarse and later moves are fine.

\subsection{The three-screen app}

\begin{enumerate}[leftmargin=1.6em,itemsep=2pt]
\item \textbf{Receipt.}  Live SAR per body region (W), current
rate (Mbps), comfort cost.  Pure information; no optimiser.  The
trust-mode user lives here.
\item \textbf{Capacity coach.}  Avatar showing the suggested local
pose move (e.g.\ ``rotate left shoulder \SI{5}{\degree} forward'').
The phone vibrates with intensity proportional to $\|\nabla R\|$
as the user moves toward the suggestion.  The suggestion updates
as UL pilots refine the twin.  This is the differentiable +
calibrated mode.
\item \textbf{Settings.}  Slider for comfort weight $\lambda_M$,
opt-in for receipt aggregation, on/off for capacity coach.  Single
user only; no operator-side cooperation among users.
\end{enumerate}

The literal-RIS-but-personal framing carries: the user's body is
\emph{their own} reflective intelligent surface, controllable only
by them, optimised for their session.  The single-user scope keeps
the privacy story simple: only the accept/reject bit and the
opt-in receipt aggregation leave the device.

\subsection{Why differentiability is load-bearing for this UX}

Without on-device gradients, the loop has only the option of
testing arbitrary new poses one at a time and waiting for UL
pilots to confirm or refute each guess.  This is a
random-search-with-measurement, slow and frustrating for the user.
With on-device gradients, the suggested move is informed by the
current calibrated twin's local linearisation; the next UL-pilot
measurement validates and refines this, but the user does not wait
for it.  Differentiability turns the closed loop from
sample-and-test into local steering: this is the operational
counterpart to the differentiability theorem of
\cref{prop:smooth}.

% ============================================================
\section{Empirical PoC: line-of-sight un-shadowing}\label{sec:eval-poc}
% ============================================================

\subsection{Geometry and parameters}

The PoC at \texttt{poc/s1\_window/} measures a single-DOF (torso
yaw) sweep in a binding-regime body-shadowed LOS scenario at
$f_c=\SI{28}{GHz}$, $M=64$ URA, BS at $\rr_{\mathrm{BS}}=(0,0,8)$~m,
user torso centroid at $\rr_{\mathrm{u}}=(30,0,1)$~m, phone offset
$(0.55,0,0.10)$~m in body-local coordinates (arm extended forward,
chest-low).  Scene loss \SI{35}{dB} (window + interior wall +
multipath floor) puts the operating SINR at the binding edge of the
MCS27 cap.  Visibility computed via trimesh ray-mesh
intersection of 49 sub-rays sampling the first Fresnel zone.

\subsection{Results}

\begin{table}[h]
\centering\small
\caption{Hero numbers per comfort budget, S1 PoC.}
\label{tab:poc}
\begin{tabular}{@{}lrrrrr@{}}
\toprule
$C\le$ & equiv.\ deg.\ & $\Delta$ rate (dB) & $\Delta$ SINR (dB)
       & $\Delta$ dose (dB) & $\Delta$ rate (\%) \\
\midrule
1   & $\sim 5$  & $+0.22$ & $+0.97$ & $+0.00$ & $+5.4$ \\
4   & $\sim 10$ & $+0.22$ & $+0.97$ & $+0.00$ & $+5.4$ \\
\textbf{9}   & \textbf{$\sim 15$} & \textbf{$+0.58$} & \textbf{$+2.63$}
            & \textbf{$+0.08$} & \textbf{$+14.5$} \\
16  & $\sim 20$ & $+0.72$ & $+3.36$ & $+0.13$ & $+18.0$ \\
36  & $\sim 30$ & $+0.91$ & $+4.66$ & $+0.21$ & $+23.1$ \\
\bottomrule
\end{tabular}
\end{table}

\noindent A \SI{15}{\degree} torso turn buys \SI{14.5}{\percent}
of rate (\SI{+87}{Mbps}, \SI{+2.6}{dB} SINR) at no
operationally-relevant exposure cost.  Dose change is
$\le\SI{1}{\percent}$ in absolute terms across the swept
configurations.

\subsection{What the PoC tests of the v4 spine}

The S1 PoC is the LOS un-shadowing limit of the cascaded channel,
in which $\partial\hLOS/\partial\theta_i$ dominates and the
body-mediated $\partial\hbody/\partial\theta_i$ contribution is
small relative to the un-shadowed direct path.  In this limit the
LOS visibility $V(\rr;\rr_{\mathrm{BS}})$ for the direct path is
the controlling quantity and the PoC reduces to a clean visibility
sweep.  This is robust to all the new physics of \cref{sec:body-to-phone}
because the body-mediated term is a small correction.  The next
PoC (NLOS, S2 in \texttt{scenarios.md}) is the spine-test for the
Kirchhoff render proper.

\subsection{Roadmap to spine-tests of \cref{sec:body-to-phone}}

\begin{enumerate}[leftmargin=1.6em,itemsep=2pt]
\item Re-run S1 with the calibrated twin (Tier C
$\hat{\bm\gamma}^{(\mathrm{C})}$ from synthetic UL pilots): verify
that the cascaded $\hh(\thetab)$ prediction matches a full-wave
reference within \SI{0.5}{dB} after calibration.
\item Implement S2 NLOS: BS path occluded, only specular bounces
available.  Expected: $\partial\hLOS/\partial\theta_i\equiv 0$,
$\partial\hbody/\partial\theta_i$ dominates; pose moves re-orient
the body to provide a useful reflection.  Spine-test for
\cref{eq:hbody} as a load-bearing primitive.
\item Implement S3 multi-user co-operative: bodies of two users
co-shape each other's cascaded channels.  Spine-test for the
multi-user gradient in \cref{sec:diff}.
\item Move from torch.no\_grad() in the SMPL-X mesh layer to
full autograd through the chain
$\thetab\to\{\rr_t\}\to\hbody\to R$; verify gradients agree with
finite differences within \SI{1}{\percent}.
\end{enumerate}

% ============================================================
\section{Connection to v5 and what we do not relitigate}\label{sec:v5}
% ============================================================

This paper is the UE-side counterpart to v5's BS-side compliance
paper, on the same per-triangle Fresnel-PO machinery.  v5 owns the
BS-side compliance loop ($\Qabs$ + per-slot primal projection of
regularised ZF + the dual-ascent QCQP solver for the binding
regime).  This paper adds the on-device twin, the
$\bm\gamma$-calibrated body-mediated channel
\cref{eq:hbody}, the UE-anchored Kirchhoff render, the SAR
receipt, and the pose-gradient suggestion.  Neither side competes
with the other.

What we do not relitigate from v5: the dictionary-side calibration
$\bm\beta$ (only re-stated as a parallel to $\bm\gamma$); the
per-slot primal projection of ZF; the dual-ascent solver for the
binding regime; the tier-A/B/C/D telemetry classification; the ISAC
bystander detection link budget.

% ============================================================
\section{What this is not}\label{sec:notthis}
% ============================================================

\begin{enumerate}[leftmargin=1.6em,itemsep=3pt]
\item Not a compliance paper.  v5 owns BS-side compliance; the
receipt of \cref{sec:saridx} is informational.
\item Not a back-hemisphere scattering paper.  The body-to-UE
channel \cref{eq:hbody} is a one-pixel render to the phone, not a
total-scattered-power Gram with a single far-field $\kout$.  The
back-hemisphere objects ($\Qref$, $\Frib$) of prior work are
correct for monostatic radar / ISAC bystander imaging but are the
wrong primitive here.
\item Not a beamforming paper.  The precoder is v5's projected ZF.
\item Not a pose-estimation paper.  The pose estimator is the
on-device IMU + AR pipeline.
\item Not a JCAS paper.  Pose tracking is on-device, not from
backscatter.
\item Not a regulatory measurement.  The SAR receipt is
model-based; the basic-restriction-pathway audit is what the
regulator certifies.
\item Not new hardware.  Vanilla 5G FR2 panel + vanilla smartphone
with IMU and AR.
\item Not a literal RIS paper.  The body cannot tune
electronically; the analogy is geometric.
\end{enumerate}

% ============================================================
\section{Loose threads and conjectures}\label{sec:threads}
% ============================================================

\begin{enumerate}[leftmargin=1.6em,itemsep=3pt]
\item \emph{Multi-bounce body re-illumination.}  Drop the
single-bounce approximation: a Neumann series in the body's
albedo $1-T_0$ gives the correction.  Convergence is guaranteed by
$\|\mathcal{R}\Qref\|\le\bar\alpha\cdot V_{\mathrm{self}}$ with
self-visibility $V_{\mathrm{self}}\ll 1$ for typical poses.
Magnitude $\sim\SI{10}{\percent}$ correction on $\Pabs$.
\item \emph{Multi-user RIHB.}  Joint loss
$\mathcal{L}(\thetab_1,\thetab_2,\xvec)$.  Conjecture: a
Stackelberg-style lift exists when both users follow gradients.
Scenario S3 measures.
\item \emph{Per-region $\bm\gamma$ stability.}  How rapidly does
$\hat{\bm\gamma}^{(\mathrm{C})}$ change with pose?  If it is
quasi-static (the body's response per region is dominated by
geometry not by surface-roughness phase variation), the calibration
overhead can be amortised over many slots.  Empirical question.
\item \emph{Sub-\SI{6}{\GHz}.}  Pseudo-Brewster collapses at
mmWave; below \SI{6}{GHz} the body's $T_0$ is polarisation-dependent
and the body-RIS analogy weakens.  A separate development for the
sub-6 carrier is needed.
\item \emph{Comfort metric calibration.}  Cost weights $c_i$
calibrated against an ergonomic study (RULA / REBA); separate
companion study.
\item \emph{Privacy preservation of the receipt aggregation.}
Differential-privacy treatments of the per-slot receipt;
secure-aggregation protocols.  Out of scope.
\item \emph{Body re-routing increases own dose.}  Re-routing
multipath through one's own body via RIHB pose moves probably
\emph{increases} that user's own absorbed dose.  Interesting
discussion observation; does not feed back into the optimisation
because exposure is informational only.
\item \emph{Pose-conditional dose-effect studies.}  Aggregated
receipts at population scale support exposure-dose-response
epidemiology that the field has been unable to run.  Study-design
question its own paper.
\item \emph{Adversarial poses.}  Worst-case pose on the comfort
ball gives the BS a regret bound on rate-loss under
non-cooperative pose.  An interesting characterisation of v5's
robustness to its own pose-estimator quality.
\item \emph{Federated $\bm\beta_{\mathrm{body}}$.}  Operator
estimates population $\bar{\bm\beta}_{\mathrm{body}}$ via federated
averaging without seeing any individual shape; whether the
smoothed-mannequin shape is operationally useful is unknown.
\item \emph{Tier B mode dimension.}  How many SVD modes does the
body-to-UE Kirchhoff operator require to capture
\SI{99}{\percent} of its spectral mass at typical
geometries?  Conjecture: $\sim 10$--$30$ at \SI{28}{GHz}, scaling
sublinearly with body solid angle from the phone.  Empirical
measurement is a one-afternoon experiment.
\item \emph{The two-paper option.}  Part~1 (SAR receipt) and
Part~2 (RIHB capacity) may merit separate papers in different
venues.  The unified spine here keeps both contributions on one
piece of machinery; splitting is mechanical.
\end{enumerate}

% ============================================================
\section{Summary}\label{sec:sum}
% ============================================================

A user with a smartphone is, at mmWave, a slowly-tunable RIS, and
a real-time exposure data source.  The right physics for both is
not a back-hemisphere scattered-power Gram (which counts total
power leaving the body) but a UE-anchored Kirchhoff radiation
integral (which gives the field at the phone antenna directly).
The integral is per-triangle, per-Fresnel-coefficient, and
naturally handles the spatially-varying direction-of-arrival
$\etahat(\rr;\rr_p)$ that distinguishes a phone-on-body geometry
from a far-field reception scenario.  Uplink pilots, which are
already in the PHY frame, calibrate the on-device twin via a
ladder of body-side scaling parameters $\bm\gamma$, mirroring v5's
dictionary-side $\bm\beta$ calibration.  The closed loop has four
orders of magnitude of bandwidth headroom relative to the data
rate it negotiates over.  The PoC measures the LOS un-shadowing
limit and finds \SI{14.5}{\percent} of rate at \SI{15}{\degree} of
torso yaw, with no operationally-relevant exposure cost.  Next is
the NLOS spine-test that exercises the Kirchhoff render proper.

% ============================================================
\appendix
% ============================================================

% ============================================================
\section{Calibration recovery: empirical evidence}\label{sec:eval-cal}
% ============================================================

This appendix reports the empirical comparison of the three
parametrisations of \cref{sec:gamma-choice} on synthetic UL
pilots at known per-region perturbation $\bm\gamma_{\mathrm{true}}
\in\Complex^6$.  Two geometries are tested:
\begin{itemize}[leftmargin=1.6em,itemsep=2pt]
\item \emph{Far-BS (\SI{30}{\m})}: BS at \SI{8}{\m} height, body
\SI{30}{\m} away.  $\Lambda$ has $K_{99}=1$.  The body subtends
$\sim 2.5^\circ$ from BS, well below the array's
$\sim 4^\circ$ beamwidth.
\item \emph{Close-BS (\SI{5}{\m})}: BS at \SI{3}{\m} height, body
\SI{5}{\m} away.  $\Lambda$ has $K_{99}=4$.  The body subtends
$\sim 13^\circ$ from BS, above beamwidth.
\end{itemize}
\Cref{fig:cal} (companion code, Fig.~3) shows the relative
$\hbody$-reconstruction error vs.\ pilot SNR for each
parametrisation, averaged over 30 trials.  The empirical winners:
the single-scalar reaches $\SI{8}{\percent}$ floor in either
geometry; per-region reaches $\SI{5}{\percent}$ at far-BS and
$\SI{3}{\percent}$ at close-BS; SVD with $K=8$ reaches
$\SI{1}{\percent}$ at far-BS \SI{40}{\dB} SNR but overfits at
$K=32$ in low-SNR.  The principled choice $K=K_{99}$ tracks the
geometry: $K=1$ at far-BS (matches single-scalar performance),
$K=4$ at close-BS (matches per-region).

% ============================================================
\section{Below 6 GHz: layered Fabry--P\'erot transmission}\label{sec:layered}
% ============================================================

The pseudo-Brewster collapse of \cref{lem:pBrewster} holds at
mmWave because skin tissue's refractive index has $|\tilde n|\sim 5$,
making $T_s\approx T_p$ to within \SI{5}{\percent}.  Below
\SI{6}{\GHz} this collapse weakens because (i) the polarisation
sensitivity of the planar interface returns ($T_s\ne T_p$ at
modest incidence angles), and (ii) the EM penetration depth
exceeds the subcutaneous fat layer thickness, so the planar
half-space approximation breaks: the fat layer sits between a
skin layer and a muscle half-space and forms a Fabry--P\'erot
resonator.  The right per-triangle replacement for $T_0$ is the
layered transmission
\begin{equation}\label{eq:T-lay}
  T_{\mathrm{lay}}(f,d_{\mathrm{SF}})
  = 1-|\tilde\Gamma_1(f,d_{\mathrm{SF}})|^2,
\end{equation}
where $\tilde\Gamma_1$ is the generalised Fresnel reflection
coefficient at the air-skin interface, computed from the
fat-muscle interface upward by the Chew recursion
\cite{Chew1995}.  $T_{\mathrm{lay}}$ depends on subcutaneous fat
thickness $d_{\mathrm{SF}}$ and exhibits a constructive resonance
near \SI{0.9}{\GHz} (quarter-wave matching) and destructive
suppression near \SI{3}{\GHz} (anti-matching) for the canonical
$\SI{2}{\milli\m}$ skin / $\SI{10}{\milli\m}$ fat / muscle stack.
The fat-thickness-weighted population average
$\langle T_{\mathrm{lay}}\rangle_{d_{\mathrm{SF}}}$ smooths the
oscillation and recovers the dosimetry-literature
plateaus
\cite{Flintoft2014,Zhang2017thesis}.  For the per-triangle
Kirchhoff render of \cref{eq:hbody} the modification is local:
substitute $T_0$ with $T_{\mathrm{lay}}(f,d_{\mathrm{SF}}^{(t)})$
per triangle, using a per-region $d_{\mathrm{SF}}$ map (which
SMPL-X can supply by anatomical region).  See
\cite{WydaegheTAP2025} for the full Chew recursion, the
layer-by-layer cube integral, and validation against \num{168}
volunteers from the dosimetry literature.

% ============================================================
\section{Compliance metrics: SARwb, psSAR10g, and peak local APD}\label{sec:compliance-metrics}
% ============================================================

The receipt of \cref{sec:saridx} reports $\Pabs$, the whole-body
absorbed power.  The ICNIRP 2020 basic restrictions
\cite{ICNIRP2020} are stated in three forms, each derivable from
$\Sab(\rr)$:

\begin{itemize}[leftmargin=1.6em,itemsep=2pt]
\item \emph{Whole-body SAR.}  $\mathrm{SAR}_{\mathrm{wb}}=
\Pabs/m_{\mathrm{body}}$; basic restriction
\SI{0.4}{\W\per\kg} for occupational and \SI{0.08}{\W\per\kg} for
general public, time-averaged over the 6-min window.  Direct
output of the receipt; the only metric that uses $\Pabs$ alone.
\item \emph{Peak spatial-averaged 10\,g SAR ($\mathrm{psSAR}_{10\mathrm{g}}$).}
The maximum of $\Pabs|_C/m_C$ over a sliding \SI{10}{\g}
contiguous tissue cube $C$.  Basic restriction
\SI{2}{\W\per\kg} for general public (head/trunk).  Requires
sub-surface SAR per voxel; the receipt's per-triangle $\Sab(\rr)$
combined with the standard skin-depth integral
$\int_0^{\delta}\sigma|E|^2\diff z$ gives a sub-surface volumetric
SAR map per voxel; the cube-averaging is a standard convolution
with the IEEE/IEC 63195 \SI{10}{\g}-cube template.
\item \emph{Peak local absorbed power density (peak local APD).}
Above \SI{6}{\GHz}, the IEC/IEEE 63195 standard replaces the
\SI{10}{\g}-cube SAR with a peak \SI{4}{\square\centi\meter}
surface-averaged absorbed power density:
$\mathrm{APD}^{\max}_{4\,\mathrm{cm}^2}=\max_{\Sigma'\subset
\Sigma,\,|\Sigma'|=4\,\mathrm{cm}^2}\frac{1}{|\Sigma'|}
\int_{\Sigma'}\Sab(\rr)\diff A$.
Basic restriction $\SI{20}{\W\per\square\meter}$ general public.
Direct surface average of the per-triangle $\Sab(\rr)$ on a sliding
\SI{4}{\square\centi\meter} window.
\end{itemize}

In operational mmWave downlink the
\emph{whole-body} restriction is the binding one (see
\cref{sec:diff,sec:app}); the receipt displays whichever metric
the user prefers, defaulting to the
$\Pabs/L_{\mathrm{RL}}$ fraction of the local reference level
budget.  The peak local APD is the most relevant for
above-\SI{6}{\GHz} compliance audits and is what the regulator
queries when the operator's compliance pipeline is independently
verified; v5 supplies the BS-side enforcement, this paper supplies
the on-device display.

% ============================================================
\section{Notation summary}\label{sec:notation}
% ============================================================

\begin{center}\small
\begin{tabular}{@{}ll@{}}
\toprule
Symbol & Meaning \\
\midrule
$\Sigma(\thetab,\betab)$ & SMPL-X body surface, pose $\thetab$,
shape $\betab$ \\
$\rr_t,\nhat_t,A_t$ & per-triangle position, normal, area \\
$\khat_n,\bm\psi_n,j(n)$ & path direction, polarisation, source
element \\
$T_0(\theta)$ & pseudo-Brewster power transmission \\
$r_0(\theta)e^{i\phi_r}$ & pseudo-Brewster amplitude reflection \\
$\Sinc(\rr),\Sab(\rr)$ & incident, absorbed power density \\
$\Qabs(\thetab)\in\Complex^{M\times M}$ & per-body absorbed-power
operator \\
$\hbody(\thetab,\rr_p)\in\Complex^M$ & body-mediated cascaded
channel \\
$\hLOS(\rr_p)\in\Complex^M$ & direct (non-body) channel \\
$\hh(\thetab,\rr_p)$ & cascaded channel
$\hLOS+\hbody$ \\
$\etahat(\rr;\rr_p)$ & per-point direction-of-arrival at phone \\
$\bm g_{\mathrm{UE}}(\hat\xi)$ & UE antenna pattern \\
$\bm\beta\in\Complex^N$ & dictionary-side calibration \\
$\bm\gamma$ & body-side calibration (Tier C/B/A) \\
$\xvec\in\Complex^M$ & BS precoder \\
$\sigma_\eta^2$ & receiver noise power \\
$\Pabs(\thetab,\xvec)=\xvec^H\Qabs\xvec$ & whole-body absorbed
power \\
$R(\thetab,\xvec)$ & user rate, capped at $S_{\max}=\SI{7.4}{bps/Hz}$ \\
$C(\thetab-\thetab^{(0)})$ & comfort cost in pose deviation \\
$L_{\mathrm{RL}}$ & ICNIRP reference level (informational) \\
$T_0\approx 0.54$ at \SI{28}{GHz}, skin & pseudo-Brewster value \\
\bottomrule
\end{tabular}
\end{center}

\end{document}
